\section{Question 2: Classification Under Extreme Imbalance}

\subsection{Problem Setup}

For Question 2, I used the credit-card fraud dataset stored at \texttt{data/external/creditcard.csv}, following the unbalanced-classification setting described by Dal Pozzolo et al. \cite{dalpozzolo2015calibrating}. The task is binary fraud detection with target column \texttt{Class}, where the positive class corresponds to fraudulent transactions. The dataset is extremely imbalanced, with a positive rate of 0.030569 and an imbalance ratio of approximately 31.71:1 in favor of the negative class. I kept a single stratified 20\% holdout split for final evaluation and restricted all model selection, resampling, calibration, and threshold tuning to the training portion only.

Leakage prevention was handled explicitly. The final test set was never touched during model comparison. For each cross-validation fold, resampling was performed only inside an \texttt{imblearn} pipeline after preprocessing, so synthetic or undersampled observations were generated from training-fold data only. I scaled \texttt{Time} and \texttt{Amount} while leaving the PCA-style \texttt{V1}--\texttt{V28} features as passthrough numeric inputs.

\subsection{Models and Sampling Strategies}

I compared four classifier families: Logistic Regression, Random Forest, XGBoost \cite{chen2016xgboost}, and a feed-forward MLP classifier. Each model was evaluated under five imbalance-handling strategies: no resampling, SMOTE \cite{chawla2002smote}, ADASYN \cite{he2008adasyn}, random undersampling, and cost-sensitive learning. Cost-sensitive learning used \texttt{class\_weight="balanced"} for Logistic Regression and Random Forest, \texttt{scale\_pos\_weight} for XGBoost, and balanced sample weights for the MLP.

Model ranking used 5-fold stratified cross-validation on the training split. Candidates were ordered primarily by PR-AUC, then MCC, then recall, since PR-AUC is more informative than ROC-AUC under severe imbalance and the downstream objective is recall-first fraud detection. The top two uncalibrated candidates were XGBoost Classifier + No Resampling and XGBoost Classifier + Cost Sensitive.

\begin{table}[htbp]
  \centering
  \caption{Top cross-validation results for Question 2, ordered by PR-AUC.}
  \label{tab:q2-cv}
  \begin{tabular}{lrrrrr}
    \toprule
    Candidate & PR-AUC & MCC & Recall & Precision & Balanced Acc. \\
    \midrule
    XGBoost Classifier + No Resampling & 0.9137 $\pm$ 0.0245 & 0.9088 & 0.8467 & 0.9818 & 0.9231 \\
    XGBoost Classifier + Cost Sensitive & 0.9104 $\pm$ 0.0230 & 0.9040 & 0.8625 & 0.9540 & 0.9306 \\
    \bottomrule
  \end{tabular}
\end{table}

\subsection{Evaluation and Calibration}

All candidate models were evaluated using precision, recall, macro F1, micro F1, ROC-AUC, PR-AUC, MCC, and balanced accuracy. After cross-validation, I recalibrated the best two candidates using both Platt scaling (sigmoid) and isotonic regression on a calibration-only validation slice drawn from the training partition. The better calibration method for each candidate was selected by Brier score.

The best calibrated variant overall was XGBoost Classifier + Cost Sensitive with isotonic calibration, achieving a calibration-set Brier score of 0.003168. Threshold selection was then performed on the calibration split using two rules: the F1-max threshold and a recall-first cost-sensitive threshold with false-negative cost 5 and false-positive cost 1. Because missing fraud is more costly than over-flagging a legitimate transaction, I used the cost-sensitive threshold as the default operating point.

\begin{table}[htbp]
  \centering
  \caption{Threshold comparison for the selected final model on the holdout split.}
  \label{tab:q2-thresholds}
  \begin{tabular}{lrrrrr}
    \toprule
    Policy & Threshold & Precision & Recall & False Negatives & False Positives \\
    \midrule
    Recall-first cost-sensitive & 0.5000 & 0.9811 & 0.9123 & 5 & 1 \\
    F1-max & 0.5000 & 0.9811 & 0.9123 & 5 & 1 \\
    \bottomrule
  \end{tabular}
\end{table}

The final selected model was XGBoost Classifier + Cost Sensitive, evaluated on the untouched holdout set. Under the default recall-first threshold, it achieved precision 0.9268, recall 0.8000, PR-AUC 0.8558, MCC 0.8571, balanced accuracy 0.8990, and Brier score 0.007574. Figures~\ref{fig:q2-pr-roc} and \ref{fig:q2-calibration-confusion} visualize the final decision behavior.

\begin{figure}[htbp]
  \centering
  \begin{minipage}[t]{0.48\textwidth}
    \centering
    \includegraphics[width=\textwidth]{../outputs/q2_classification/20260420-135405_credit_card_fraud_baseline/figures/best_model_precision_recall_curve.png}
  \end{minipage}
  \hfill
  \begin{minipage}[t]{0.48\textwidth}
    \centering
    \includegraphics[width=\textwidth]{../outputs/q2_classification/20260420-135405_credit_card_fraud_baseline/figures/best_model_roc_curve.png}
  \end{minipage}
  \caption{Precision-recall and ROC curves for the final calibrated Question 2 model on the holdout set.}
  \label{fig:q2-pr-roc}
\end{figure}

\begin{figure}[htbp]
  \centering
  \begin{minipage}[t]{0.48\textwidth}
    \centering
    \includegraphics[width=\textwidth]{../outputs/q2_classification/20260420-135405_credit_card_fraud_baseline/figures/best_model_reliability.png}
  \end{minipage}
  \hfill
  \begin{minipage}[t]{0.48\textwidth}
    \centering
    \includegraphics[width=\textwidth]{../outputs/q2_classification/20260420-135405_credit_card_fraud_baseline/figures/best_model_confusion_matrix.png}
  \end{minipage}
  \caption{Reliability diagram and confusion matrix for the final calibrated Question 2 model at the recall-first operating threshold.}
  \label{fig:q2-calibration-confusion}
\end{figure}

\subsection{Discussion}

The results show why PR-AUC and recall-aware thresholding matter for extreme imbalance. A model can look acceptable under ROC-AUC while still failing to identify enough fraud cases at a usable operating threshold. The final pipeline therefore emphasized PR-AUC during model ranking and then used a cost-sensitive threshold to explicitly favor recall. This choice reflects the application setting: missing a fraudulent transaction is typically more costly than reviewing an extra flagged transaction.

Calibration was also necessary because raw classifier scores are not always well aligned with true fraud probabilities. By comparing sigmoid and isotonic calibration on a dedicated validation slice, the final workflow separates ranking quality from probability quality. This improves both the reliability diagram and the downstream threshold decision. In practice, the most meaningful trade-off is between recall and operational alert volume, so the final report should interpret false positives as review workload and false negatives as direct fraud risk.
